% Expository adaptation, October 7, 2026. Original mathematics: OpenAI, September 23, 2026.
% Prose revised and explanations expanded; see README.md for the source and scope.
\section{Reducing transition entries and accumulating the reward}
\label{sec:hierarchy}

This section constructs the exact target of the estimator. It is useful to keep its final guarantee in view: at each stage, the largest transition entry decreases by a factor $2^H$, while the number of vertices incident to positive entries increases by at most two. The residual transition mass then tends to zero geometrically. Propositions~\ref{prop:hier-copy-step} and~\ref{prop:hier-reward-target} state the one-stage guarantee and its consequence for acceptance probabilities.

There are three operations to justify. A correction transfers transition contributions into the reward. A detector restricts that correction in columns with too many substantial incoming entries. A copy step divides the remaining large entries while controlling how many copies become active. We begin with the identity that makes the first operation exact.

Given a transition matrix $C$ and a nonnegative correction $E$, put
\[
 D=C-E+EC.
\]
Then $I-D=(I+E)(I-C)$. Thus a relation $W=(I-C)p$ becomes
$(I+E)W=(I-D)p$: the same vector $p$ is represented by the new transition
and the updated reward. The correction must make $D$ nonnegative and keep
its row sums at most one. It must also reduce entries sufficiently fast
that, after a controlled increase in the number of active vertices, the
remaining transition contribution becomes small.  The construction uses
only the model reduction, with no one-sided acceptance promise.

At one stage the input is a finite ordered set $V$, a rate $0<f\le1$
on the grid of Section~\ref{sec:env}, and a nonnegative matrix $C$ with
\begin{equation}
 C(x,y)=0\quad\text{unless }x<y,\qquad
 C\mathbf1\le\mathbf1,\qquad C(x,y)\le f.
 \label{eq:hier-input}
\end{equation}
All matrix and vector inequalities are entrywise. Write
\[
 \mathcal R(C)=\{x:\text{some }C(x,y)>0\},\quad
 \mathcal T(C)=\{y:\text{some }C(x,y)>0\},\quad
 \mathcal A(C)=\mathcal R(C)\cup\mathcal T(C).
\]
The stage will reduce the entry bound from $f$ to $f/2^H$ and increase
$|\mathcal A(C)|$ by at most a factor of two. We first analyze the
correction with an arbitrary fixed column multiplier. We then choose that
multiplier by a finite detector calculation, and use copies to divide the
entries in columns where correction is not fully enabled.

\subsection{Correction with a fixed column multiplier}

The correction will remove most of a sufficiently large entry, but it must leave a small positive remainder. That remainder has two uses. It ensures nonnegativity of the new transition matrix, and its positive size bounds the number of corrected positions in a row. The cutoff $Q$ expresses this rule continuously, so that replacing an input by a nearby estimate has a controlled effect. The gap between its two thresholds leaves room for later errors.

Put $q=2^H$.  Fix rational numbers
$0<\theta_1<\theta_2<q^{-1}/32$ and a nondecreasing Lipschitz function
$Q:[0,\infty)\to[0,\infty)$ satisfying
\begin{equation}
 (t-\theta_2)_+\leq Q(t)\leq(t-\theta_1)_+.
 \label{eq:hier-Q}
\end{equation}
Its eighth root is required to be Lipschitz on bounded intervals. This extra regularity is used when errors in support decisions are bounded by eighth moments. It is a choice we can build into the cutoff, rather than an additional hypothesis on the input matrix.
Here and below the scalar cutoffs can be chosen piecewise polynomial with
rational coefficients.  To see this explicitly, let $\beta$ be the
normalized integral of $t^7(1-t)^7$ on $[0,1]$, extended by zero to the
left and one to the right.  Both $\beta^{1/8}$ and
$(1-\beta)^{1/8}$ are Lipschitz: near their vanishing endpoints the
corresponding functions vanish to order eight, and away from those
endpoints the assertion follows by differentiation.  For example,
\[
 Q(t)=(t-\theta_1)_+
       \beta\left(\frac{t-\theta_1}{\theta_2-\theta_1}\right)
\]
has all the stated properties.  Near $\theta_1$ its eighth root is of
order $(t-\theta_1)_+^{9/8}$; elsewhere its root has bounded derivative
on every bounded interval.  Let $L_Q$ be a fixed Lipschitz constant for
$Q$.  For real arguments, extend $Q$ by $Q(t)=0$ when $t<0$.
This extension preserves its Lipschitz and root-Lipschitz properties and
permits evaluation on real approximations.  The exact hierarchy itself
uses only nonnegative arguments.

For the moment, suppose we have already decided how much correction each column permits. Write this decision as a fixed multiplier $g:V\to[0,1]$. A value of one enables the full correction, and a value of zero disables it. Starting from zero, define
\begin{equation}
 E_0=0,\qquad G'_i=C+E_{i-1}C,\qquad
 E_i(x,y)=fQ(G'_i(x,y)/f)g(y)\quad(i\ge1).
 \label{eq:hier-fixed-g-iteration}
\end{equation}
The cutoff leaves a positive residual at each corrected entry.  Keeping
$g$ fixed makes the update map monotone.  We next show that a fixed
number of iterations suffices for every such multiplier, independently
of the size of $V$.  At a nonnegative fixed point $E_*$ of this recurrence, put
$T=C-E_*+E_*C$.  In a column with $g(y)=1$, the lower bound on $Q$
gives $T(x,y)\le\theta_2f$.  The truncation estimate below will transfer
this entry bound to a finite correction, with a small additional error.
The fixed point is only a comparison object; the hierarchy will use a
finite iteration.

\begin{lemma}[Fixed point and factorial truncation]
\label{lem:hier-fixed-point}
For every fixed column multiplier $0\leq g\leq1$, there is a unique
strictly forward nonnegative solution of
\begin{equation}
 E_*(x,y)=fQ((C+E_*C)(x,y)/f)g(y).
 \label{eq:hier-fixed-point}
\end{equation}
For the recurrence in Equation~\eqref{eq:hier-fixed-g-iteration}, $0\leq E_i\leq E_*$, and for every $k\geq0$,
\begin{equation}
 \max_{x,y}\frac{E_*(x,y)-E_k(x,y)}f
 \leq
 e^{1/\theta_1}\frac{(L_Q/\theta_1)^k}{k!}.
 \label{eq:hier-factorial}
\end{equation}
In particular $k_*$ can be chosen as a fixed integer so that the
right side of Equation~\eqref{eq:hier-factorial} at $k=k_*$ is at most
$q^{-1}/32$, uniformly in $C,f,V$, and $g$.
\end{lemma}

\begin{proof}
We prove the lemma in three steps: construct the comparison fixed point, bound the number of corrected positions in one row, and count the increasing paths along which a truncation error can propagate.

Order the vertices compatibly with the strictly forward relation.  For a
fixed row $x$, the entry $(E_*C)(x,y)$ involves only $E_*(x,z)$ with
$z<y$.  Equation~\eqref{eq:hier-fixed-point} therefore defines the row
successively, in increasing column order.  Columns at or before $x$
receive zero.  This proves existence and uniqueness in the stated class,
without a contraction hypothesis.  Monotonicity of the update map proves
$E_i\leq E_*$ and the nonnegativity of the differences.

To bound the support of a fixed row, use the positive remainder deliberately left by $Q$. Set $T=C-E_*+E_*C$. If $a=(C+E_*C)(x,y)$, then
Equation~\eqref{eq:hier-Q} implies
\[
 0\leq E_*(x,y)\leq(a-\theta_1f)_+.
\]
Thus $T\geq0$, and at every positive position of $E_*$ one has
$T(x,y)\geq\theta_1f$.  Also
\begin{equation}
 T\mathbf 1=C\mathbf 1-E_*(\mathbf 1-C\mathbf 1)
       \leq C\mathbf 1\leq\mathbf 1.
 \label{eq:hier-fixed-residual-mass}
\end{equation}
Consequently, for each fixed row $x$, its set
$A_x=\{y:E_*(x,y)>0\}$ has cardinality
$m_x\leq1/(\theta_1f)$.

We next bound an entry of the fixed point by expanding its recurrence. The useful count is the size of the corrected support of this one row, not the total size of the graph. Unroll the inequality $E_*\leq C+E_*C$ along this row's support.
Each resulting term for $E_*(x,y)$ is a product of entries of $C$
along a sequence
\[
 x<z_1<\cdots<z_j<y,
 \qquad z_1,\ldots,z_j\in A_x.
\]
Every substitution keeps the source coordinate equal to $x$: the
recursive factor is $E_*(x,z)$, while only its terminal coordinate
changes.  All intermediate vertices therefore belong to the same set
$A_x$.  Strict ordering makes a choice of intermediate vertices specify
at most one such sequence.
There are at most $\binom{m_x}{j}$ such sequences, and each product
is at most $f^{j+1}$.  The expansion terminates because every sequence
is strictly increasing.  Therefore
\begin{equation}
 E_*(x,y)\leq
 f\sum_{j=0}^{m_x}\binom{m_x}{j}f^j
 =f(1+f)^{m_x}\leq fe^{1/\theta_1}.
 \label{eq:hier-fixed-entry}
\end{equation}
Only the support of the single fixed row $x$ was used here; the number
of vertices in the whole graph does not enter the exponent.

The same path count now controls the error after finitely many iterations. Each iteration forces an error contribution to traverse one more strictly increasing step; choosing those steps produces the factorial in the denominator. Let $\Delta_i=E_*-E_i$. These matrices are nonnegative, and their
row-$x$ supports are contained in $A_x$.  The Lipschitz bound gives
\[
 \Delta_i(x,y)\leq
 L_Q\sum_{z<y}\Delta_{i-1}(x,z)C(z,y).
\]
Outside $A_x$ the left side is zero, so when this inequality is unrolled
it suffices to use intermediate vertices from $A_x$ at every step.
After $k$ steps there are at most $\binom{m_x}{k}$ strictly increasing
sequences of intermediate vertices.  Their initial differences are at
most $fe^{1/\theta_1}$ by Equation~\eqref{eq:hier-fixed-entry}, and
their $k$ transition factors are at most $f^k$.  Thus
\[
 \Delta_k(x,y)
 \leq fe^{1/\theta_1}L_Q^kf^k\binom{m_x}{k}
 \leq fe^{1/\theta_1}\frac{(L_Q/\theta_1)^k}{k!}.
\]
For $k>m_x$ the first upper bound is zero.  This proves
Equation~\eqref{eq:hier-factorial}.  The factorial eventually dominates
the fixed exponential factor, giving the stated fixed choice of $k_*$.
\end{proof}

The iteration count $k_*$ is now fixed independently of the detectors.
We next choose $g$ so that correction is supported on columns with only
boundedly many substantial entries.  A pilot selects this multiplier;
the actual recurrence then uses that one multiplier for all $k_*$ steps.

\subsection{Choosing which columns permit correction}

The fixed-multiplier argument did not explain how to choose $g$. We need correction to be enabled only in columns with a controlled number of substantial entries. A preliminary sequence of corrections, called the pilot, estimates how crowded each column is. After that sequence, its final choice of $g$ is frozen and the actual correction starts again from zero. The two sequences serve different purposes: the pilot chooses the multiplier; the actual sequence is the monotone recurrence analyzed above.

At this stage these are exact mathematical definitions. Every detector is a true expectation of an already defined input. Computing an approximation to the detector comes later, so an approximate value never defines its own exact target.

For the detectors we use independent channels $X,Y$ from
Section~\ref{sec:env}.  Denote their common retention probability at rate
$f$ by $\pi_f$.  The consequences of Lemma~\ref{lem:env-ranks} needed here
can be written, with fixed positive constants, as
\begin{align}
 c_\pi f\leq\pi_f\leq C_\pi f,\qquad
 &\Pr\{Y_f(w)=1\mid Y_f(y)=1\}\leq C_{\rm pair} f
       &&(w\ne y), \label{eq:hier-sampling-pair}\\
 &\Pr\left\{\sum_{x\in A}X_f(x)<\tfrac12|A|\pi_f\right\}
       \leq\frac{C_{\rm var}}{|A|\pi_f}
       &&(\varnothing\ne A\subseteq V\text{ deterministic}).
       \label{eq:hier-sampling-tail}
\end{align}
These constants are uniform in the stage, the admissible moduli, and common
fixed padding suffixes.  No independence between different rows of a
detector trial will be assumed.

All threshold choices and ranks below are fixed constants.  The pilot
uses the iteration count from Lemma~\ref{lem:hier-fixed-point}.  Choose
\[
 0<\delta_0<\delta_1<\theta_1/100,
 \qquad 0<\eta_-<\eta_+<\delta_0/3.
\]
If other row scores require still smaller detector thresholds, make those
choices before fixing the ranks.  Let $\phi$ be a nondecreasing cutoff
which is zero through $\delta_0$ and one from $\delta_1$ on.  Let $\rho$
be a nonincreasing cutoff which is one through $\eta_-$ and zero from
$\eta_+$ on.  Finally, let $\gamma:[0,1]\to[0,1]$ be nonincreasing,
one on $[0,1/2]$, and zero on $[4/5,1]$.  These functions and their
complements, when used, have the root regularity just described.

A detector looks for many retained rows that contribute substantially to the same column. It must first exclude rows with too many sampled entries, because those rows will not support a bounded-access representation. The row gate $R_i$ performs this exclusion; the $J$th-largest score then tests whether enough eligible rows remain.

For a finite multiset of nonnegative numbers, $\operatorname{rank}_j$
means its $j$-th largest member, with enough additional zeros that this is
always defined.  Starting from $\widehat E_0=0$, perform $k_*$ pilot
steps.  At step $i$, first set
\begin{equation}
 G_i=C+\widehat E_{i-1}C.
 \label{eq:hier-pilot-input}
\end{equation}
For a draw of the terminal channel, set
\begin{equation}
 R_i(x)=\rho\left(\operatorname{rank}_K
       \{G_i(x,w)/f:Y_f(w)=1\}\right).
 \label{eq:hier-pilot-row-gate}
\end{equation}
Define the deterministic column number
\begin{equation}
 U_i(y)=\mathbb E\left[
    \operatorname{rank}_J
       \{\phi(G_i(x,y)/f)R_i(x):X_f(x)=1\}
       \,\middle|\,Y_f(y)=1\right].
 \label{eq:hier-detector}
\end{equation}
Then put
\begin{equation}
 \widehat g_i(y)=\gamma\left(\max_{1\leq r\leq i}U_r(y)\right),
 \qquad
 \widehat E_i(x,y)=fQ(G_i(x,y)/f)\widehat g_i(y).
 \label{eq:hier-pilot-correction}
\end{equation}
The order of these definitions is essential: $G_i$ is already
deterministic when its detector is formed, and $U_i$ is defined before
$\widehat E_i$.  The running maximum makes $\widehat g_i$ nonincreasing.
Its final value is therefore no larger than any earlier pilot multiplier;
monotonicity of $Q$ will let us bound the actual corrections by the pilot
corrections.

After the pilot, set $g=\widehat g_{k_*}$ and specialize the fixed-multiplier
recurrence of Equation~\eqref{eq:hier-fixed-g-iteration} to the actual correction
\begin{equation}
 E_0=0,\qquad G'_i=C+E_{i-1}C,\qquad
 E_i(x,y)=fQ(G'_i(x,y)/f)g(y)\quad(1\leq i\leq k_*).
 \label{eq:hier-actual-iteration}
\end{equation}
Write
\begin{equation}
 E=E_{k_*},\qquad D=C-E+EC.
 \label{eq:hier-residual-definition}
\end{equation}

\begin{lemma}[Elementary pilot bounds]
\label{lem:hier-pilot-bounds}
All pilot and actual inputs and corrections are nonnegative and strictly
forward.  For $1\leq i\leq k_*$,
\begin{align*}
 G_i\mathbf 1,\ \widehat E_i\mathbf 1,
     \ G'_i\mathbf 1,\ E_i\mathbf 1&\leq i\mathbf 1,\\
 G_i(x,y),\ \widehat E_i(x,y),
     \ G'_i(x,y),\ E_i(x,y)&\leq if.
\end{align*}
Moreover,
\[
 E_{i-1}\leq E_i\leq\widehat E_i,
 \qquad G'_i\leq G_i.
\]
The row and column supports of every matrix in this statement are
contained in $\mathcal R(C)$ and $\mathcal T(C)$, respectively.
\end{lemma}

\begin{proof}
Equation~\eqref{eq:hier-Q} gives $Q(0)=0$ and
$0\leq Q(t)\leq t$ for $t\ge0$.
Consequently $\widehat E_i\leq G_i$.  If
$\widehat E_{i-1}\mathbf 1\leq(i-1)\mathbf 1$, then
\[
 G_i\mathbf 1=C\mathbf 1+\widehat E_{i-1}C\mathbf 1
       \leq\mathbf 1+\widehat E_{i-1}\mathbf 1
       \leq i\mathbf 1.
\]
For an individual entry,
\[
 G_i(x,y)\leq f+f\sum_z\widehat E_{i-1}(x,z)\leq if.
\]
This proves the pilot bounds by induction.

The sequence $\widehat g_i$ is nonincreasing, so
$g\leq\widehat g_i$ for every $i$.  Induction in
Equations~\eqref{eq:hier-pilot-input}--\eqref{eq:hier-actual-iteration},
using monotonicity of $Q$, proves $E_i\leq\widehat E_i$ and
$G'_i\leq G_i$.  The map
$F\mapsto fQ((C+FC)/f)g$, applied entrywise after forming the matrix
product, is monotone.  Its iteration from zero is therefore
nondecreasing.  The actual bounds follow from the pilot bounds.

A product of strictly forward matrices is strictly forward.  A zero row
of $C$ remains a zero row in all these formulas, by induction from the
zero correction.  A zero column of $C$ is a zero column of every product
of the form $FC$ and hence of every input and correction.  This proves
the remaining assertions.
\end{proof}

The next lemma supplies both directions of the count comparison. A sufficiently crowded column is detected with large score. Conversely, a positive score above the stated threshold certifies many substantial entries. The first implication limits correction support; the second later limits the number of extra copies.

\begin{lemma}[The detector bounds overloaded columns]
\label{lem:hier-column-test}
Let
\[
 D_{\rm row}=1+C_{\rm pair}k_*/\eta_-.
\]
Choose the integer $K\geq160D_{\rm row}$, increasing it if necessary
for the other finitely many row-gate estimates.  For any fixed positive
integer $J$, choose
\begin{equation}
 R_*\geq c_\pi^{-1}\max(4J,40C_{\rm var}).
 \label{eq:hier-column-constants}
\end{equation}
Then the pilot has the following properties.
\begin{enumerate}
\item If
$\#\{x:G_i(x,y)\geq\delta_1f\}>R_*/f$, then
$U_i(y)\geq19/20>9/10$.
\item If $U_i(y)>1/4$, then
\[
 \#\{x:G_i(x,y)>\delta_0f\}
       >\frac{J}{4C_\pi f}.
\]
\item For any pilot correction with multiplier
$\widehat g_i(y)>0$,
\[
 \#\{x:G_i(x,y)\geq\delta_1f\}\leq R_*/f.
\]
For every actual correction, $g(y)>0$ implies
\[
 \#\{x:G'_i(x,y)\geq\delta_1f\}\leq R_*/f.
\]
\end{enumerate}
\end{lemma}

\begin{proof}
First bound how often the row gate removes a particular row. This estimate is conditional on retaining the queried column, as the detector requires. Fix $i,y$ and condition on $Y_f(y)=1$.  In a fixed row, the set
$\{w:G_i(x,w)>\eta_-f\}$ has cardinality at most
$k_*/(\eta_-f)$ by Lemma~\ref{lem:hier-pilot-bounds}.
Its expected number of retained members is at most $D_{\rm row}$:
the specified terminal $y$, if in this set, contributes at most one,
and Equation~\eqref{eq:hier-sampling-pair} applies to every other member.
If the row gate is not full, at least $K$ such members have been retained.
Markov's inequality therefore gives
\begin{equation}
 \Pr\{R_i(x)<1\mid Y_f(y)=1\}\leq D_{\rm row}/K.
 \label{eq:hier-row-failure}
\end{equation}

For a crowded column, compare the number of sampled substantial rows with the number lost to row gates. It is enough to bound the expected number lost; the losses need not be independent. Let $A=\{x:G_i(x,y)\geq\delta_1f\}$ and $n=|A|$.
Let $N_A$ count its retained starts, and let $L_A$ count those retained
starts whose row gate is not full.  Independence of the two channels
and Equation~\eqref{eq:hier-row-failure} imply
\[
 \mathbb E[L_A\mid Y_f(y)=1]\leq n\pi_fD_{\rm row}/K.
\]
Thus
\[
 \Pr\{L_A>n\pi_f/4\mid Y_f(y)=1\}\leq1/40.
\]
The law of the start channel is unchanged by the conditioning, so
Equation~\eqref{eq:hier-sampling-tail} gives
\[
 \Pr\{N_A<n\pi_f/2\mid Y_f(y)=1\}
       \leq C_{\rm var}/(n\pi_f).
\]
If $n>R_*/f$, Equation~\eqref{eq:hier-column-constants} ensures
$n\pi_f\geq\max(4J,40C_{\rm var})$.  Outside an event of probability
at most $1/20$, there are at least $J$ retained high rows with full
gates.  Each has score exactly one, so the trial's $J$-th largest score
is one.  All trial scores lie in $[0,1]$, proving the first assertion.
This argument used no independence among the row-gate failures.

For the converse, a positive $J$th-largest score requires at least $J$ sampled eligible starts. Counting those starts bounds the mean score from above. For the second assertion let
$A_0=\{x:G_i(x,y)>\delta_0f\}$, and let $N_0$ be its number of
retained starts.  A positive score must come from $A_0$.  Pointwise the
trial is at most $\mathbf 1_{\{N_0\geq J\}}$, and hence at most
$N_0/J$.  Its conditional expectation is at most
$|A_0|\pi_f/J\leq |A_0|C_\pi f/J$.  The asserted lower bound on
$|A_0|$ follows.

Finally, a positive $\widehat g_i(y)$ implies $U_i(y)<4/5$, so the
first assertion excludes a high-entry count greater than $R_*/f$.
If $g(y)>0$, all pilot multipliers at $y$ are positive.  Apply the
same bound to every $G_i$ and then use $G'_i\leq G_i$ from
Lemma~\ref{lem:hier-pilot-bounds}.
\end{proof}

\subsection{The finite residual}

The actual construction uses the finite correction $E=E_{k_*}$, not the comparison fixed point. We therefore check that stopping the iteration preserves nonnegativity and bounded row sums, and that the error paid for stopping is small enough in fully corrected columns.

There is a further requirement for estimation. The residual is the signed expression $C-E+EC$, so its constituent terms can be larger than the final answer. We need a bound on each constituent relative to the nonnegative residual. Without it, a small relative error in a constituent could become a large error after cancellation. The domination estimate below supplies exactly this bound.

\begin{lemma}[Residual bounds and unsigned-term domination]
\label{lem:hier-residual}
For the $k_*$ chosen in Lemma~\ref{lem:hier-fixed-point}, put
\[
 K_*=k_*+1,\qquad
 \varepsilon_{\rm good}=\theta_2+q^{-1}/32<q^{-1}/16,
 \qquad c_{\rm dom}=\theta_1/(2k_*).
\]
The matrix $D$ in Equation~\eqref{eq:hier-residual-definition} is
nonnegative, strictly forward, and substochastic.  Its entries satisfy
\begin{equation}
 D(x,y)\leq K_*f,
 \qquad
 g(y)=1\ \Longrightarrow\ D(x,y)\leq\varepsilon_{\rm good}f.
 \label{eq:hier-residual-entry-bounds}
\end{equation}
Furthermore,
\begin{equation}
 D\geq c_{\rm dom}(C+E+EC).
 \label{eq:hier-domination}
\end{equation}
Thus every unsigned term in $C-E+EC$ is bounded by a fixed multiple
of $D$.  In the formulas $G_i=C+\widehat E_{i-1}C$ and
$G'_i=C+E_{i-1}C$, each unsigned term is bounded by the corresponding
input itself.  Moreover,
$\mathcal R(D)\subseteq\mathcal R(C)$ and
$\mathcal T(D)\subseteq\mathcal T(C)$.
\end{lemma}

\begin{proof}
The last correction was computed from $G'_{k_*}$, whereas the residual uses $E_{k_*}$ in its product with $C$. Separate their difference first. Set $A=G'_{k_*}$ and
$\delta=(E-E_{k_*-1})C\geq0$.  Then
\begin{equation}
 D=A-E+\delta,
 \qquad A-E\geq\min(A,\theta_1f).
 \label{eq:hier-subsolution-residual}
\end{equation}
Here and in the minimum the scalar $\theta_1f$ denotes the constant
entrywise bound.  The second inequality follows from
$E\leq(A-\theta_1f)_+$.  It proves $D\geq0$.
Since $C\mathbf 1\leq\mathbf 1$ and $E\geq0$,
\[
 D\mathbf 1=C\mathbf 1-E(\mathbf 1-C\mathbf 1)
       \leq C\mathbf 1\leq\mathbf 1.
\]
Also $D\leq C+EC$, whose entries are at most
$f+f\sum_z E(x,z)\leq(k_*+1)f$ by
Lemma~\ref{lem:hier-pilot-bounds}.  Strict forwardness and the support
claims follow from the same formulas.

Because $A(x,y)\leq k_*f$, Equation~\eqref{eq:hier-subsolution-residual}
gives $A-E\geq(\theta_1/k_*)A$.  On the other hand,
\[
 C+E+EC=A+E+\delta\leq2A+\delta.
\]
Since $\theta_1/k_*<1$, these two inequalities prove
Equation~\eqref{eq:hier-domination}.  The remaining unsigned-term
assertions are immediate from nonnegativity of the summands in the two
input formulas.

For the good-column bound use the fixed point with this same multiplier
$g$, and set $\Delta=E_*-E$.  With $T=C-E_*+E_*C$,
\[
 D=T+\Delta-\Delta C\leq T+\Delta.
\]
If $g(y)=1$, the lower bound in Equation~\eqref{eq:hier-Q} gives
$T(x,y)\leq\theta_2f$.  Equation~\eqref{eq:hier-factorial} gives
$\Delta(x,y)\leq q^{-1}f/32$.  This proves
Equation~\eqref{eq:hier-residual-entry-bounds}.
\end{proof}

The domination in Equation~\eqref{eq:hier-domination} is also an
estimation bound.  It lets us compare the three unsigned contributions
$C$, $E$, and $EC$ separately against their estimates, measuring each
error relative to the same nonnegative target $D$.  Their signs are
then combined by a triangle inequality in
Lemma~\ref{lem:est-mean-comparison}.  The estimated raw arrays may have
negative entries; positivity here is a property of the exact target.

\subsection{Divide the remaining entries between copies}

The residual already has small entries in columns with $g(y)=1$.
We distribute the remaining large entries among copies of their
destinations.  The detector incidence bound will limit the number of
parents with active extra copies, and a weighted projection will relate
the new transition to the residual.

A primary copy is always present. An extra copy receives transition mass only when an entry is still too large, and its outgoing row is enabled only when the detector permits it. The routing fractions below specify how the incoming mass is divided; the row factors specify which copies can continue the computation. Their compatibility is what makes the projection identity exact.

Choose an integer $D_0>K_*q$, and let
$\tau_0=q^{-1}/8$, $\tau_1=q^{-1}/4$.
Choose a rising cutoff $\chi$ which is zero through $\tau_0$ and one
from $\tau_1$ on.  For $0\leq j<D_0$ set
\begin{equation}
 s_0(t)=1-(1-D_0^{-1})\chi(t),\qquad
 s_j(t)=D_0^{-1}\chi(t)\quad(j>0).
 \label{eq:hier-routing-fractions}
\end{equation}
They are nonnegative, sum to one, and have Lipschitz eighth roots on
bounded ranges.  Define $U^{\max}(x)=\max_{1\leq i\leq k_*}U_i(x)$.
Let $\kappa$ be a rising cutoff, zero through $1/4$ and one from
$1/2$ on, and put
\begin{equation}
 r_0(x)=1,\qquad r_j(x)=\kappa(U^{\max}(x))\quad(j>0).
 \label{eq:hier-copy-row-factors}
\end{equation}
On $V^+=V\times\{0,\ldots,D_0-1\}$ write $x_i=(x,i)$ and set
\begin{equation}
 C^+(x_i,y_j)=r_i(x)D(x,y)s_j(D(x,y)/f).
 \label{eq:hier-copy-transition}
\end{equation}
Order the copies first by their parent and then by their copy symbol.

\begin{proposition}[One-stage density and active-count reduction]
\label{prop:hier-copy-step}
Choose the detector rank $J$ large enough that
\begin{equation}
 J\geq\frac{2C_\pi(D_0-1)k_*(k_*+1)}{\delta_0},
 \label{eq:hier-active-J}
\end{equation}
and then choose $R_*$ as in Equation~\eqref{eq:hier-column-constants}.
The matrix $C^+$ is strictly forward and substochastic, and
\[
 C^+(x_i,y_j)\leq f/q,
 \qquad |\mathcal A(C^+)|\leq2|\mathcal A(C)|.
\]
If $P^r$ is the matrix from $V^+$ to $V$ defined by
\begin{equation}
 P^r(x_i,x)=r_i(x),\qquad
 P^r(x_i,z)=0\quad(z\ne x),
 \label{eq:hier-projection}
\end{equation}
then
\begin{equation}
 C^+P^r=P^rD.
 \label{eq:hier-one-step-intertwining}
\end{equation}
\end{proposition}

\begin{proof}
We check the entry and row-sum bounds first, then prove that projection recovers the old transition, and finally count the active copies.

The row sum at $x_i$ is $r_i(x)\sum_yD(x,y)\leq1$.  Strict
forwardness follows from that of $D$.  Put $t=D(x,y)/f$.  If
$t<\tau_1$, then every routed entry is at most $ft<f/q$.  If
$t\geq\tau_1$, all fractions equal $1/D_0$, so every entry is at
most $K_*f/D_0<f/q$.  This proves the density bound.

If a nonprimary fraction $s_j(D(x,y)/f)$ is positive on a positive
entry, then $D(x,y)/f>\tau_0>\varepsilon_{\rm good}$.
Lemma~\ref{lem:hier-residual} therefore implies $g(y)<1$.
Since $\gamma$ equals one through $1/2$, this gives
$U^{\max}(y)>1/2$, and hence $r_j(y)=1$.  The primary factor is
always one.  Thus for every positive entry of $D$,
\begin{equation}
 \sum_{j=0}^{D_0-1}s_j(D(x,y)/f)r_j(y)=1.
 \label{eq:hier-routing-projects}
\end{equation}
For a zero entry no identity of its fractions is needed.  Multiplying
Equation~\eqref{eq:hier-copy-transition} by $P^r$ and using
Equation~\eqref{eq:hier-routing-projects} proves
Equation~\eqref{eq:hier-one-step-intertwining} entrywise.

The formal vertex set grew by the factor $D_0$, but most extra vertices can be inactive. We count parents of active extra copies by charging each such parent to many substantial pilot entries. Row-mass bounds then limit the total number of those incidences. Put
\[
 \mathcal O=\{y:U^{\max}(y)>1/4\},\qquad
 n'=|\mathcal A(C)|.
\]
Every $y\in\mathcal O$ has, for some pilot step $i$, more than
$J/(4C_\pi f)$ entries with $G_i(x,y)>\delta_0f$, by
Lemma~\ref{lem:hier-column-test}.  Counting these incidences over all
steps and all active rows, and using
Lemma~\ref{lem:hier-pilot-bounds}, gives
\begin{align*}
 |\mathcal O|\frac{J}{4C_\pi f}
 &\leq\sum_{i=1}^{k_*}\#\{(x,y):G_i(x,y)>\delta_0f\}\\
 &\leq\frac{n'}{\delta_0f}\sum_{i=1}^{k_*}i
   =\frac{n'k_*(k_*+1)}{2\delta_0f}.
\end{align*}
The assertion is also valid when $n'=0$, since then all matrices and
detectors vanish.  Thus
\[
 |\mathcal O|\leq
       \frac{2C_\pi n'k_*(k_*+1)}{\delta_0J}.
\]

Every active child has an active parent, by the support conclusion of
Lemma~\ref{lem:hier-residual}.  A nonprimary child active as a row has
$r_i(x)>0$, so its parent belongs to $\mathcal O$.  A nonprimary
child active as a column has a positive extra routing fraction and, as
shown above, its parent has $U^{\max}>1/2$, so again belongs to
$\mathcal O$.  Therefore
\[
 |\mathcal A(C^+)|\leq n'+(D_0-1)|\mathcal O|\leq2n',
\]
where the last inequality is Equation~\eqref{eq:hier-active-J}.
\end{proof}

\begin{remark}[Order of the one-stage constants]
\label{rem:hier-constant-order}
The bounds used to select $k_*$ in
Lemma~\ref{lem:hier-fixed-point} hold for every multiplier in $[0,1]$.
Thus $\theta_1,\theta_2,Q,k_*,K_*$ and $D_0$ can be fixed before any
detector rank.  Next fix all finitely many correction, detector, and
access-score threshold ladders, taking $\delta_0,\delta_1$ below the
correction support thresholds as required.  Choose $K$ large enough for
the corresponding row estimates.  Choose $J$ by
Equation~\eqref{eq:hier-active-J}, and then $R_*$ by
Equation~\eqref{eq:hier-column-constants}.  An incoming rank used later
can be chosen after $R_*$: its full-gate probability is controlled by
the column bound in Lemma~\ref{lem:hier-column-test}, independently of
the auxiliary row gates.  There is therefore no requirement to enlarge
$K$ again in response to that incoming rank.  Subsequent estimator
constants may depend on these already fixed ranks and thresholds.
\end{remark}

\subsection{Iterate the construction and bound the remaining error}

The copy step controls transition entries and active vertices.  We now
iterate it and transport the terminal reward through the same weighted
projections.  The resulting reward will approximate the original
acceptance probability, including probabilities strictly between zero
and one.

The remaining calculation connects the local construction to the original acceptance probability. Each copy step has a projection back to its parent. Composing these projections tells us which weighted base value a current vertex represents; composing the reward updates tells us how much of that value has already been collected.

Starting with $C_0=S$ on $V_0$, repeat the one-stage construction with
\[
 f_l=q^{-l},\qquad
 V_{l+1}=V_l\times\{0,\ldots,D_0-1\},\qquad C_{l+1}=C_l^+.
\]
Let $E_l,D_l,P_l^r$ denote the correction, residual, and projection at
stage $l$.  Define matrices from $V_l$ to $V_0$ by
\begin{equation}
 P_0^c=I,\qquad P_{l+1}^c=P_l^rP_l^c,
 \qquad B_0=I,\qquad B_{l+1}=P_l^r(I+E_l)B_l.
 \label{eq:hier-cumulative-matrices}
\end{equation}
The superscript $c$ here denotes the cumulative projection, not a
transpose.  Let $e$ be the base final-success reward and
$p_0=(I-S)^{-1}e$, so $0\leq p_0\leq1$ by the model reduction.

\begin{proposition}[Reward identity and a shallow target]
\label{prop:hier-reward-target}
For every stage,
\begin{equation}
 C_lP_l^c=P_l^c-B_l(I-S).
 \label{eq:hier-reward-algebra}
\end{equation}
The vectors $W_l=B_le$ satisfy
\begin{equation}
 0\leq W_l\leq1,\qquad
 W_l=P_l^cp_0-C_lP_l^cp_0,
 \label{eq:hier-reward-bound}
\end{equation}
and their updates are
\begin{equation}
 W_{l+1}(x_i)=r_i(x)\bigl(W_l(x)+(E_lW_l)(x)\bigr).
 \label{eq:hier-reward-update}
\end{equation}
If $\bar x_l$ is the all-primary copy of a base vertex $x$, then
\begin{equation}
 0\leq p_0(x)-W_l(\bar x_l)
       \leq |V_0|\,2^{-(H-1)l}.
 \label{eq:hier-final-error}
\end{equation}
Consequently there is $L=O(B)$ for which the last upper bound is
strictly less than $1/16$ for every base vertex $x$.
In particular, $p_0(x)=0$ implies $W_L(\bar x_L)=0$, and
$p_0(x)\ge1/2$ implies $W_L(\bar x_L)>7/16$.
Every full stage domain has size $2^{O(B)}$ and admits $O(B)$-bit
addresses.
\end{proposition}

\begin{proof}
First combine the correction identity and the copy projection into one invariant. The assertion at $l=0$ is $S=I-(I-S)$.  The residual identity can be
written as
\[
 D_l=I-(I+E_l)(I-C_l).
\]
Using Equation~\eqref{eq:hier-one-step-intertwining} and then the
inductive hypothesis,
\begin{align*}
 C_{l+1}P_{l+1}^c
 &=P_l^rD_lP_l^c\\
 &=P_l^r\bigl(P_l^c-(I+E_l)(P_l^c-C_lP_l^c)\bigr)\\
 &=P_{l+1}^c-P_l^r(I+E_l)B_l(I-S)\\
 &=P_{l+1}^c-B_{l+1}(I-S).
\end{align*}
This proves Equation~\eqref{eq:hier-reward-algebra}.

Every $B_l$ is nonnegative.  Every $P_l^r$ has at most one nonzero
entry per row, of value in $[0,1]$, and thus $P_l^c$ is
substochastic.  Apply Equation~\eqref{eq:hier-reward-algebra} to
$p_0$, using $(I-S)p_0=e$.  This gives the identity in
Equation~\eqref{eq:hier-reward-bound}.  Its lower bound follows from
$B_le\geq0$; its upper bound follows from
$W_l\leq P_l^cp_0\leq1$.  Equation~\eqref{eq:hier-reward-update}
is Equation~\eqref{eq:hier-cumulative-matrices} applied to $e$.

By Proposition~\ref{prop:hier-copy-step}, induction gives
\[
 |\mathcal A(C_l)|\leq2^l|\mathcal A(S)|\leq2^l|V_0|,
 \qquad C_l(x,y)\leq q^{-l}.
\]
There are at most $|\mathcal A(C_l)|$ columns containing any positive
entry, so each row sum of $C_l$ is at most
$|V_0|2^lq^{-l}=|V_0|2^{-(H-1)l}$.
The all-primary projection has every row factor equal to one and hence
$(P_l^cp_0)(\bar x_l)=p_0(x)$.  The identity in
Equation~\eqref{eq:hier-reward-bound}, together with
$0\leq P_l^cp_0\leq1$, now proves
Equation~\eqref{eq:hier-final-error}.

Choose an integer $L$ such that
$(H-1)L>\log_2(16|V_0|)$.  Since $\log_2|V_0|=O(B)$, this is
$O(B)$.  If $p_0(x)=0$, nonnegativity forces
$W_L(\bar x_L)=0$.  If $p_0(x)\geq1/2$, the strict error bound
gives $W_L(\bar x_L)>7/16$.  Finally,
\[
 |V_l|=|V_0|D_0^l=2^{O(B)}\qquad(0\leq l\leq L),
\]
because $D_0$ is fixed.  Parent projection deletes one radix symbol,
and appending a specified symbol is affine in the numerical ID.  Thus
these domain changes preserve the address convention required by the
sampler and by the later implementation.
\end{proof}

The remainder of the proof will estimate the finite targets just
constructed.  In particular, it need only implement the $k_*$ pilot and
$k_*$ actual updates at a stage and the reward update in
Equation~\eqref{eq:hier-reward-update}; it does not evaluate $E_*$.
The same error bound in Equation~\eqref{eq:hier-final-error} gives any
fixed inverse-polynomial accuracy by increasing the constant in
$L=O(B)$.  Section~\ref{sec:conclusion} makes that choice together with
the statistical and numerical precisions.

